\section*{References}
\addcontentsline{toc}{section}{References}
\begingroup\small\setlength{\parindent}{0pt}\raggedright
\hangindent=1.5em\hangafter=1 Andersen, T. G., and Bollerslev, T. (1998). Answering the skeptics: Yes, standard volatility models do provide accurate forecasts. \textit{International Economic Review}, 39(4), 885--905.\par\smallskip
\hangindent=1.5em\hangafter=1 Barroso, P., and Santa-Clara, P. (2015). Momentum has its moments. \textit{Journal of Financial Economics}, 116(1), 111--120.\par\smallskip
\hangindent=1.5em\hangafter=1 Black, F. (1976). Studies of stock price volatility changes. \textit{Proceedings of the 1976 Meetings of the American Statistical Association, Business and Economic Statistics Section}, 177--181.\par\smallskip
\hangindent=1.5em\hangafter=1 Bollerslev, T. (1986). Generalized autoregressive conditional heteroskedasticity. \textit{Journal of Econometrics}, 31(3), 307--327.\par\smallskip
\hangindent=1.5em\hangafter=1 Bollerslev, T., Tauchen, G., and Zhou, H. (2009). Expected stock returns and variance risk premia. \textit{Review of Financial Studies}, 22(11), 4463--4492.\par\smallskip
\hangindent=1.5em\hangafter=1 Bongaerts, D., Kang, X., and van Dijk, M. (2020). Conditional volatility targeting. \textit{Financial Analysts Journal}, 76(4), 54--71.\par\smallskip
\hangindent=1.5em\hangafter=1 Campbell, J. Y., and Thompson, S. B. (2008). Predicting excess stock returns out of sample: Can anything beat the historical average? \textit{Review of Financial Studies}, 21(4), 1509--1531.\par\smallskip
\hangindent=1.5em\hangafter=1 Clark, T. E., and West, K. D. (2007). Approximately normal tests for equal predictive accuracy in nested models. \textit{Journal of Econometrics}, 138(1), 291--311.\par\smallskip
\hangindent=1.5em\hangafter=1 Cederburg, S., O'Doherty, M. S., Wang, F., and Yan, X. S. (2020). On the performance of volatility-managed portfolios. \textit{Journal of Financial Economics}, 138(1), 95--117.\par\smallskip
\hangindent=1.5em\hangafter=1 Chicago Board Options Exchange, CBOE Volatility Index: VIX [VIXCLS], retrieved from FRED, Federal Reserve Bank of St. Louis. \url{https://fred.stlouisfed.org/series/VIXCLS}\par\smallskip
\hangindent=1.5em\hangafter=1 Elmachtoub, A. N., and Grigas, P. (2022). Smart ``predict, then optimize''. \textit{Management Science}, 68(1), 9--26.\par\smallskip
\hangindent=1.5em\hangafter=1 French, K. R. Kenneth R. French Data Library. \url{https://mba.tuck.dartmouth.edu/pages/faculty/ken.french/data_library.html}\par\smallskip
\hangindent=1.5em\hangafter=1 French, K. R., Schwert, G. W., and Stambaugh, R. F. (1987). Expected stock returns and volatility. \textit{Journal of Financial Economics}, 19(1), 3--29.\par\smallskip
\hangindent=1.5em\hangafter=1 Glosten, L. R., Jagannathan, R., and Runkle, D. E. (1993). On the relation between the expected value and the volatility of the nominal excess return on stocks. \textit{Journal of Finance}, 48(5), 1779--1801.\par\smallskip
\hangindent=1.5em\hangafter=1 Harvey, C. R., Hoyle, E., Korgaonkar, R., Rattray, S., Sargaison, M., and Van Hemert, O. (2018). The impact of volatility targeting. \textit{Journal of Portfolio Management}, 45(1), 14--33.\par\smallskip
\hangindent=1.5em\hangafter=1 Künsch, H. R. (1989). The jackknife and the bootstrap for general stationary observations. \textit{Annals of Statistics}, 17(3), 1217--1241.\par\smallskip
\hangindent=1.5em\hangafter=1 Ledoit, O., and Wolf, M. (2008). Robust performance hypothesis testing with the Sharpe ratio. \textit{Journal of Empirical Finance}, 15(5), 850--859.\par\smallskip
\hangindent=1.5em\hangafter=1 Moreira, A., and Muir, T. (2017). Volatility-managed portfolios. \textit{Journal of Finance}, 72(4), 1611--1644.\par\smallskip
\hangindent=1.5em\hangafter=1 Moskowitz, T. J., Ooi, Y. H., and Pedersen, L. H. (2012). Time series momentum. \textit{Journal of Financial Economics}, 104(2), 228--250.\par\smallskip
\hangindent=1.5em\hangafter=1 Newey, W. K., and West, K. D. (1987). A simple, positive semi-definite, heteroskedasticity and autocorrelation consistent covariance matrix. \textit{Econometrica}, 55(3), 703--708.\par\smallskip
\hangindent=1.5em\hangafter=1 Patton, A. J. (2011). Volatility forecast comparison using imperfect volatility proxies. \textit{Journal of Econometrics}, 160(1), 246--256.\par\smallskip
\hangindent=1.5em\hangafter=1 Politis, D. N., and Romano, J. P. (1992). A circular block-resampling procedure for stationary data. In R. LePage and L. Billard (Eds.), \textit{Exploring the Limits of Bootstrap} (pp. 263--270). Wiley.\par\smallskip
\hangindent=1.5em\hangafter=1 Rapach, D. E., Strauss, J. K., and Zhou, G. (2010). Out-of-sample equity premium prediction: Combination forecasts and links to the real economy. \textit{Review of Financial Studies}, 23(2), 821--862.\par\smallskip
\hangindent=1.5em\hangafter=1 Romano, J. P., and Wolf, M. (2005). Stepwise multiple testing as formalized data snooping. \textit{Econometrica}, 73(4), 1237--1282.\par\smallskip
\hangindent=1.5em\hangafter=1 Shumway, T. (1997). The delisting bias in CRSP data. \textit{Journal of Finance}, 52(1), 327--340.\par\smallskip
\hangindent=1.5em\hangafter=1 Welch, I., and Goyal, A. (2008). A comprehensive look at the empirical performance of equity premium prediction. \textit{Review of Financial Studies}, 21(4), 1455--1508.\par\smallskip
\hangindent=1.5em\hangafter=1 Xu, X. (2024). Improving volatility-managed portfolios in real time. Working paper, ESSCA School of Management. SSRN 4778937. \url{https://ssrn.com/abstract=4778937}\par\smallskip
\hangindent=1.5em\hangafter=1 Zhao, P., Zhu, H., Ng, W. S. H., and Lee, D. L. (2024). From GARCH to neural network for volatility forecast. \textit{Proceedings of the AAAI Conference on Artificial Intelligence}, 38(15), 16998--17006.\par\smallskip
\endgroup
